% !TeX root = ../概率论与数理统计笔记.tex
% 由主文档通过 \input 载入；请编译主文档。
\setcounter{chapter}{7} % 保留原书章号。
\chapter{生成函数与特征函数}
\label{chap:ps-generating-functions}

本章研究如何把分布编码为函数，以便处理矩、独立和与极限。
先修内容为期望、幂级数和独立性；特征函数部分另需复数与欧拉公式。首次学习可先掌握基本计算，再阅读唯一性与收敛理论。

\section{概率生成函数}
\label{sec:ps-pgf}
对非负整数值随机变量定义\emph{概率生成函数}（probability generating function, PGF）。
说明在适当收敛区域内由系数恢复概率、由导数求阶乘矩，以及独立和对应函数乘积。
边界处求导需单独检查所需矩是否有限。

\section{矩母函数}
\label{sec:ps-mgf}
定义\emph{矩母函数}（moment-generating function, MGF）$M_X(t)=\mathbb E[e^{tX}]$，其中 $t$ 为实数。
区分在单个点有限与在零点的某个开邻域内有限；在后者条件下讨论求导取矩及分布唯一性。
以具有有限若干矩但矩母函数不存在于零点邻域的分布说明方法边界。

\section{特征函数及其基本性质}
\label{sec:ps-characteristic-function}
令 $i$ 为虚数单位，定义\emph{特征函数}（characteristic function）$\varphi_X(t)=\mathbb E[e^{itX}]$。
说明它对所有实数 $t$ 都存在，推导平移、缩放和独立和的公式。
陈述特征函数的唯一性定理；连续性定理作为第\ref{chap:ps-limits}章分布收敛证明的选读工具，明确极限函数在零点连续的条件。

\section{例题、习题与学习检查}
规划伯努利、泊松、正态与伽马分布的变换计算，以及独立和分布的识别。
习题比较“所有矩存在”“矩母函数在零点邻域有限”“矩能唯一确定分布”这些不同命题，不将它们混同。
学习检查要求说明每次求导与期望交换的依据，并选择最适合当前分布的工具。






